\section{Introduction}
\label{sec:intro}

Selecting a base policy to post-train for a specific deployment scenario is, in current
practice, largely a leaderboard-reading exercise: rank candidates on nuScenes, NAVSIM, or
the CARLA Leaderboard, and pick the winner. This is convenient because public scores are
free to look up and directly comparable \emph{within} one benchmark. It is also
epistemically thin for a different question -- ``which policy will behave best on
\emph{my} target scenario after post-training'' -- because a benchmark score is a single
scalar aggregate over a fixed distribution of driving situations, while a deployment
scenario is typically rare, safety-critical, and out of that distribution. A policy that
tops a benchmark by being reliably mediocre everywhere can lose to a policy that is
worse on average but happens to have the right internal machinery for the scenario that
matters.

The obvious response -- run both candidates on the target scenario and compare -- is
exactly what full closed-loop re-benchmarking requires, and it is precisely the cost this
line of work tries to avoid: closed-loop simulators (CARLA, Bench2Drive) are expensive to
build scenario content for, and a numeric score, even a target-domain one, still only
tells us \emph{how much} a policy degrades, not \emph{where in the model} the
degradation originates or \emph{what it would cost to fix}. Two policies can obtain
identical target-domain scores for entirely different internal reasons, and the correct
post-training prescription is different in each case.

This paper asks a narrower, answerable version of the question: instead of predicting a
full ranking, can we build a small set of literature-grounded, causally validated axes
that \emph{localize} where a benchmark-to-deployment gap enters a policy's causal chain,
using only a modest amount of paired out-of-distribution data and no closed-loop
simulation? We propose four such axes -- \textbf{G}rounding, \textbf{F}aithfulness,
\textbf{I}nvariance, and \textbf{C}oncentration (GFIC) -- each anchored to an independent
body of causal-inference or interpretability literature rather than invented ad hoc for
one scenario:

\begin{itemize}
    \item \textbf{G (Grounding)}, after shortcut learning~\cite{geirhos2020shortcut}: does
    the policy's internal representation encode objective, task-independent scene
    structure -- what is object versus background -- beyond what its architecture's
    prior alone provides, rather than relying on a task-irrelevant shortcut?
    \item \textbf{F (Faithfulness)}, after causal confusion in imitation
    learning~\cite{dehaan2019causal}: does the hazard entity actually drive the action,
    or does removing it from the input at test time leave the action unchanged?
    \item \textbf{I (Invariance)}, after invariant risk minimization~\cite{arjovsky2019irm}:
    is the policy's behaviour stable under a controlled appearance-domain shift, and does
    its representational stability agree with its behavioural stability?
    \item \textbf{C (Concentration)}, after localization-for-editing
    work~\cite{hase2023localization,meng2022rome}: when the policy does degrade under
    domain shift, is the causal responsibility concentrated in a small, identifiable set
    of layers, or diffuse?
\end{itemize}

We instantiate this protocol on a ghost-probe (occluded pedestrian sudden emergence)
target scenario mined from nuScenes and validate it on \textbf{six} end-to-end driving
policies chosen to separate three confounds a smaller pool cannot: encoder family,
action-head family, and native training domain. Three policies (DiffusionDrive, LTF,
DiffusionDriveV2) share one encoder under three different action heads; two (SimLingo,
and the two VLAs Alpamayo-R1 and AutoVLA) are architecturally unrelated to that family and
to each other. This design is deliberate: a diagnostic protocol's portability failures,
and the boundary between ``the encoder determines it'' and ``the action head determines
it,'' are invisible in a pool of two.

As \cref{fig:pipeline} summarizes, GFIC is positioned as an explicit, causally validated
step between a policy finishing training and a deployment decision -- one that today is
typically skipped in favour of reading a leaderboard.

\begin{figure}[t]
\centering
\begin{tikzpicture}[
  node distance=5mm and 7mm,
  every node/.style={font=\small},
  box/.style={draw, rounded corners, align=center, minimum height=8mm, minimum width=24mm, inner sep=2pt},
  hl/.style={box, fill=cvprblue!12, draw=cvprblue, thick},
  small/.style={draw, rounded corners, align=center, minimum height=7mm, minimum width=24mm, font=\scriptsize, inner sep=2pt},
  hlsmall/.style={small, fill=cvprblue!12, draw=cvprblue, thick},
  arr/.style={-{Latex[length=2mm]}, thick}
]
\node[box] (train) {Trained /\\post-trained policy};
\node[hl, right=of train] (diag) {GFIC diagnosis\\\textbf{(this paper)}};
\node[box, right=of diag] (deploy) {Deployment\\decision};
\draw[arr] (train) -- (diag);
\draw[arr] (diag) -- (deploy);

\node[small, below=13mm of train] (pub) {Public\\leaderboard rank};
\node[hlsmall, below=13mm of diag] (axis) {GFIC axis\\rank};
\node[small, below=13mm of deploy] (real) {Real deployment\\rank (target-specific)};
\draw[-{Latex[length=1.5mm]}, gray] (train.south) .. controls +(0,-6mm) and +(0,6mm) .. (pub.north);
\draw[-{Latex[length=1.5mm]}, cvprblue] (diag.south) -- (axis.north);
\draw[-{Latex[length=1.5mm]}, gray] (deploy.south) .. controls +(0,-6mm) and +(0,6mm) .. (real.north);

\node[font=\scriptsize\bfseries] at ($(pub.east)!0.5!(axis.west)$) {$\overset{?}{=}$};
\node[font=\scriptsize\bfseries] at ($(axis.east)!0.5!(real.west)$) {$\overset{?}{=}$};
\end{tikzpicture}
\caption{GFIC as an explicit diagnostic step between training and deployment. A public
leaderboard rank, a GFIC axis profile, and a real target-scenario deployment rank need
not agree with one another (\cref{sec:experiments} shows cases of each disagreement);
this paper measures the middle term causally, from a small amount of target-domain
out-of-distribution data and with no closed-loop simulator.}
\label{fig:pipeline}
\end{figure}

Our contributions are:

\begin{enumerate}
    \item A four-axis diagnostic protocol (GFIC) that is domain- and scenario-general by
    construction -- the ghost-probe scenario is one instantiation, not the definition --
    each axis anchored to an independent, established literature and each measured with
    an explicit causal-validity check rather than a correlational proxy alone. We test
    portability along two independent axes of generality: a second, causally distinct
    scenario type (an already-tracked lead vehicle braking hard, rather than a new entity
    entering the path) and a second, independently collected data source (NAVSIM's
    underlying OpenScene/nuPlan logs). Concentration's architectural law replicates
    completely (14 of 14 cells across both tests); Grounding and Faithfulness -- G-VS and
    F-3, our objective-ground-truth readouts that require no semantic judgement of danger
    -- are markedly less portable: G-VS's point estimate is stable across both tests but
    its three-state verdict is not, and F-3 is only testable where a candidate shows any
    baseline response to a hazard frame at all, a structural limit of a necessity test
    rather than a defect of one corpus. The four axes are not equally portable, and no
    single-corpus, single-scenario study could have shown that.
    \item A demonstration that public benchmark scores can fail to even define a ranking
    across a pool of six candidates (different benchmark, domain, protocol, and metric),
    which motivates the need for an orthogonal, target-domain diagnostic rather than a
    benchmark-substitute score.
    \item Empirical evidence that a single axis (Concentration) can prescribe genuinely
    different repair sites for policies public scores cannot distinguish, and generalizes
    into a pre-registered, held-out-confirmed law separating two architecture families
    with no sign overlap; and that a second axis (Invariance) orders two policies sharing
    an identical encoder oppositely depending on whether it is read from representations
    or from behaviour -- two independent arguments against compressing diagnosis into one
    number.
    \item Two objective-ground-truth readouts for Grounding and Faithfulness --
    segmentation-probe selectivity against Segment Anything pseudo-ground-truth (G-VS) and
    input-level occlusion necessity (F-3) -- that need no judgement of ``what counts as
    dangerous'' and, for the first time, apply uniformly across all six candidates
    regardless of architecture or language interface. Without any hazard criterion, they
    identify a joint dissociation on DiffusionDrive (the representation demonstrably
    encodes object-level information that the action does not act on) together with an
    analogous blind-action signature on LTF, and both dissociations replicate across the
    scenario- and data-source-generality tests of contribution 1.
    \item A diagnosis-guided pilot post-training intervention with an attribution-control
    arm, reported honestly as a partial success: the targeted causal link is repaired,
    but the repair does not translate into a behavioural gain at the tested budget.
\end{enumerate}
